\subsection{A local obstruction to an inner square root}

\begin{lemma}\label{lem:2d2-neutral-indicators}
Let $\mathcal B$ be pseudo-unitary with simple objects $1,g,r,gr$
and rules
\[
 g^2=1,\qquad gr=rg,\qquad r^2=1+2r+2gr.
\]
Then $\nu_2(g)=-1$, $\nu_2(r)=\nu_2(gr)=1$, and the cubic
rotations on $\operatorname{Hom}(1,r^3)$ and
$\operatorname{Hom}(1,(gr)^3)$ are identities.
\end{lemma}
\begin{proof}
Put $d=2+\sqrt5$ and $D=4+8d=20+8\sqrt5$.
The formal codegrees are $D,D',4,4$, and
$1/D-1/D'=-1/\sqrt5$. The proof of Lemma~\ref{lem:q17-conductor-equality}, with the
primitive character modulo five in place of that modulo seventeen,
gives the same equality conclusion at conductor five: every central simple with nonrational dimension
ratio $z/D$ has twist of order dividing five.

We use the signed version of the saturated indicator identity.
Let $\ell(z)$ be a rational coordinate of the nontrivial quadratic
part of $z/D$, scaled so that its values on $(1,g,r,gr)$ are
$(2,2,-1,-1)$. For a cyclotomic automorphism $\sigma$ with residue
$k$ modulo five, write $\pi$ and $\varepsilon_A$ for its modular
permutation and signs. The Chinese remainder theorem lifts each chosen
unit residue to an exponent coprime to the full cyclotomic conductor.
The dimension transformation gives
\[
 z_{\pi A}=\varepsilon_A\frac D{\sigma D}\sigma(z_A),
 \qquad \ell(z_{\pi A})=\varepsilon_A\chi_5(k)\ell(z_A).
\]
The first induction identity, weighted by $\ell$, is
\[
 \sum_A\ell(z_A)z_A\theta_A=D\ell(1).
\]
The signs in the two preceding transformations cancel, giving
\[
 \sigma\bigl(\ell(z_A)z_A\bigr)
 =\chi_5(k)\frac{\sigma D}{D}\ell(z_{\pi A})z_{\pi A}.
\]
Apply $\sigma$ to the weighted identity and reindex by $\pi$.
If $h$ is the inverse of $k$ modulo five, its nonzero terms obey
$\sigma(\theta_A)=\theta_{\pi A}^{h}$, since their twists divide
five and $\theta_{\pi A}=\sigma^2(\theta_A)$. The terms belonging
to $I(1)$ have twist one and satisfy the same equation.
Cancel $\sigma D$ and use the induction formula for the $h$th
indicator. Inversion permutes the units modulo five and
$\chi_5(h)=\chi_5(k)$, proving
\[
 \sum_i\ell(\dim X_i)\nu_k(X_i)=\chi_5(k)\ell(1)
 \qquad (k,5)=1.
\]
At $k=2$ this reads
$2+2\nu_2(g)-\nu_2(r)-\nu_2(gr)=-2$.
All three unknowns are signs, forcing the stated values.
At $k=3$ it reads $\nu_3(r)+\nu_3(gr)=4$.
Each rotation acts on a two-dimensional space and has order
three. Equality in the trace bound forces both rotations to be
identities.
\end{proof}

\begin{lemma}\label{lem:2d2-no-inner-square}
For such a pseudo-unitary $\mathcal B$, no tensor autoequivalence
fixing all four simple labels has square tensor-isomorphic to
$\operatorname{Ad}(g)$.
\end{lemma}
\begin{proof}
We use only the indicated local pivotal identities from
\cite[\S3.1.1 and pp.~539--540]{EIP}, interpreted on opposite
composition-dual Hom spaces using \cite[Lemma~3.13(ii),
Proposition~3.22 and Corollary~4.9]{FGPartialDual}.
No Hermitian form or unitarity assumption is made.
Put
\[
 H^+=\operatorname{Hom}(r^2,r),\quad
 J^+=\operatorname{Hom}(r^2,gr),\qquad
 H^-=\operatorname{Hom}(r,r^2),\quad
 J^-=\operatorname{Hom}(gr,r^2).
\]
Each space has dimension two. Normalize the $g$--$r$ crossing
$c$ so that its pivotal one-click rotation is its inverse.
Its bend parameter is $\mu=\nu_2(gr)/\nu_2(r)=1$.

Here is the paired-space normalization, including the distinction
between transpose and inverse transpose. For $F=\operatorname{Ad}(g)$,
take
$q_r=(1_r\otimes\operatorname{ev}_g)(c\otimes1_g)$ and
$q_{gr}=1_g\otimes q_r$. On opposite Hom spaces its outgoing and
incoming actions use respectively these identifications and their
inverses. Suppressing the tensorator and its inverse, they are
\[
 T^+(f)=q_BF(f)(q_A^{-1}\otimes q_A^{-1}),\qquad
 T^-(h)=(q_A\otimes q_A)F(h)q_B^{-1}.
\]
Composition pairing is invariant, so the incoming loop
matrix is $K^-=(K^+)^{-\mathsf t}$. Transposes of the object maps
are relative to the canonical pivotal structure of $F$; the unit
map $q_1=\operatorname{ev}_g$ has form defect one.
The two defects are $s_r=\mu$ and $s_{gr}=\mu\nu_2(g)$:
the first contraction is the defining two-crossing $\mu$ diagram,
and the second reverses one further antisymmetric $g$ form.
For a right bend $R$ on $\operatorname{Hom}(A^2,B)$ this gives
$RT^+=(s_B/s_A^2)T^-R$.

The direct loop contractions give $K_H^2=-1$, $K_J^2=1$.
The nondegenerate symmetric right theta pairings consequently pair
opposite loop eigenlines. We may choose composition-dual bases with
\[
 \begin{gathered}
 K_H^+=\operatorname{diag}(i,-i),\quad
 K_H^-=\operatorname{diag}(-i,i),\\
 K_J^+=K_J^-=\operatorname{diag}(1,-1),\qquad
 R_H^+=R_H^-=R_J^+=R_J^-=E=\begin{pmatrix}0&1\\1&0\end{pmatrix}.
 \end{gathered}
\]
Their remaining basis changes are
$\operatorname{diag}(u,u^{-1})$ and
$\operatorname{diag}(v,v^{-1})$.
The cubic rotation on $H$ is the identity by
Lemma~\ref{lem:2d2-neutral-indicators}. The two decorated cubic
phases on $J$ remain arbitrary third roots throughout the argument.

In these bases the tetrahedral slots of \cite[pp.~528--530]{EIP}
have the required forms
\[
 A=\begin{pmatrix}
 t&0&0&A_1\\0&t&t&0\\0&t&t&0\\a_1&0&0&t
 \end{pmatrix},\qquad
 D_0=\begin{pmatrix}
 x&0&0&Z\\0&X&y&0\\0&Y&-X&0\\z&0&0&-x
 \end{pmatrix}.
\]
The odd partial-dual slot map is
$(i,j,k,l)\mapsto(k,1-j,i,1-l)$, pairing
$a_1\leftrightarrow A_1$, $x\leftrightarrow X$,
$y\leftrightarrow Y$ and $z\leftrightarrow Z$.
These are independent dual coefficients, not complex conjugates.
For completeness, the even slot maps are
\[
 c(i,j,k,l)=(j,1-k,1-i,l),\qquad
 q(i,j,k,l)=(1-k,1-l,1-i,1-j).
\]
The $A$ coefficients obey both even symmetries. For $D_0$, only
its eight odd-parity zeros and the two relations
$(D_0)_{1111}=-(D_0)_{0000}$ and
$(D_0)_{1010}=-(D_0)_{0101}$ are needed. In both flips $i=k$, so the relevant decorated phase
cancels as its own cube. No equality of decorated phases is used.

The fusion-ring corner on the four-dimensional annular module
over $r$ has block matrices
\[
 P=\begin{pmatrix}0&I\\I&0\end{pmatrix},\qquad
 Q=\begin{pmatrix}\phi&\phi_0\\\phi_0&\phi\end{pmatrix},
 \qquad \phi=2tI.
\]
The FP character is absent, since its associated central simple
is the central unit, whose forgetful image contains no $r$
\cite[\S2.3]{OstrikPiv3}. Thus $\phi+\phi_0=(2-\sqrt5)I$.
On $P=-1$ the fusion relation is $Q^2=I$, giving
\[
 t=\frac{2-\sqrt5+\tau}{4},\qquad \tau=\pm1.
\]
Three diagonal composition-dual resolutions on \cite[p.~539]{EIP}
now suffice:
\[
 t^2+xX=\frac{3-\sqrt5}{2},\qquad
 xX+zZ=\frac12,\qquad xX+yY=\frac12.
\]
They come from the first resolution at $(0,0,0,0)$ and the third
at $(0,0,0,0)$ and $(0,1,0,1)$. In the first identity the
unit-channel contribution is $1/d=\sqrt5-2$, so its uncombined
form is $1=2t^2+2xX+(\sqrt5-2)$. Each diagonal contraction pairs
every vertex phase with its reciprocal. Therefore
\[
 yY=zZ=
 \begin{cases}(\sqrt5-1)/8&\tau=1,\\
 (3\sqrt5-5)/8&\tau=-1.
 \end{cases}
\]
Both are nonzero, so $y,z\ne0$ in every branch.

Let $T$ fix the four simple labels. It preserves the canonical
pivotal structure. Normalize its object identifications to preserve
the self-dual forms. The crossing is taken to $\lambda c$ with
$\lambda=\pm1$, and its $H$ action satisfies
$T_HK_H=\lambda K_HT_H$. If $\lambda=-1$, its matrix is
antidiagonal, so $T_H^2$ is scalar. If $\lambda=1$, the normalized
crossing and hence both loop eigenspace pairs are fixed. The right
pairings then give diagonal actions
$\operatorname{diag}(u,u^{-1})$ and
$\operatorname{diag}(v,v^{-1})$. The two nonzero coefficients
transform with weights $u^{-2}v^2$ and $u^{-2}v^{-2}$.
Their invariance forces $u^4=1$, so again $T_H^2$ is scalar.

The projective action on $H$ depends only on the tensor-isomorphism
class of the autoequivalence and respects composition: changing
the identification of a simple object changes this action only
by a scalar. But $\operatorname{Ad}(g)$ acts projectively by
$K_H=\operatorname{diag}(i,-i)$, which is not scalar.
It therefore cannot be the square of $T$.
\end{proof}

\begin{theorem}\label{thm:r24-inner-square}
The ring $R8\_2\_24$ has no complex fusion categorification.
\end{theorem}
\begin{proof}
The source has exactly two FP dimensions, $1$ and $2+\sqrt5$,
so pass to a pseudo-unitary Galois conjugate using
\cite[Lemma~3.1]{SchNear}. The source tensor has a faithful
$C_2$ grading with neutral labels $\{0,5,6,7\}$ and the ring of
Lemma~\ref{lem:2d2-neutral-indicators}, with $r=X_5$ and $g=X_7$.
The odd invertible $a=X_1$ has $a^2=g$. Since the source ring is
commutative, $\operatorname{Ad}(a)$ fixes every neutral simple
label. Its square is tensor-isomorphic to
$\operatorname{Ad}(g)$, contradicting
Lemma~\ref{lem:2d2-no-inner-square}.
\end{proof}

The source-bound replay in
\texttt{verify\_r24\_projective\_inner\_square.py} checks
the neutral tensor, complete four-character table, saturation arithmetic,
coefficient-slot substitutions and both nonzero scalar branches.
The categorical argument above supplies the obstruction; the finite
replay does not assume a classification of all $2D2$ associators.
